\documentclass[11pt]{article}

\usepackage[margin=1in]{geometry}
\usepackage{amsmath,amssymb}
\usepackage{mathtools}
\usepackage{booktabs}
\usepackage{enumitem}
\usepackage[hidelinks]{hyperref}

\title{Residual Gauge Evolution for Background-Subtracted Z4c}
\author{AthenaK TDE Debugging Note}
\date{\today}

\newcommand{\ubg}{u_{\mathrm{bg}}}
\newcommand{\ures}{u_{\mathrm{res}}}
\newcommand{\ufull}{u_{\mathrm{full}}}
\newcommand{\abg}{\alpha_{\mathrm{bg}}}
\newcommand{\ares}{\alpha_{\mathrm{res}}}
\newcommand{\afull}{\alpha_{\mathrm{full}}}
\newcommand{\bbg}{\beta^i_{\mathrm{bg}}}
\newcommand{\bres}{\beta^i_{\mathrm{res}}}
\newcommand{\bfull}{\beta^i_{\mathrm{full}}}
\newcommand{\rhs}{\mathcal{R}}
\newcommand{\gaugea}{\mathcal{G}_{\alpha}}
\newcommand{\gaugeb}{\mathcal{G}_{\beta^i}}

\begin{document}
\maketitle

\section{Purpose}

The goal of the background-subtracted Z4c system is to evolve a compact
secondary object on top of an analytically known primary black-hole background
without allowing truncation error from the primary background to dominate the
secondary's self-gravity.  The current debugging question is whether the
gauge variables, especially lapse and shift, can be treated in a way that is
both internally consistent and stable.

The immediate concern is the following.  If the background lapse and shift are
prescribed analytically, but the star or secondary black hole requires a
different local gauge, does a local residual gauge evolution effectively
``change the background'' and spoil the subtraction?  The answer proposed here
is no, provided the background remains a fixed reference state and the gauge
equations are written for the residual variables in a background-subtracted
form.  The residual gauge then changes only the full evolved spacetime, not
the analytic background used for subtraction.

\section{Residual Variables}

Let the Z4c evolved variables be denoted schematically by \(u\).  The split is
\begin{equation}
  \ufull(x,t) = \ubg(x,t) + \ures(x,t).
\end{equation}
Here \(\ubg\) is a prescribed analytic reference solution, for example
Kerr-Schild Schwarzschild in the black-hole rest frame.  It is not evolved
numerically.  It is recomputed from analytic formulae or otherwise prescribed
by the background callback.  The residual \(\ures\) is the only dynamical
quantity stored and updated by the residual formulation.

The residual evolution should be
\begin{equation}
  \partial_t \ures
  =
  \rhs[\ufull] - \rhs[\ubg],
  \label{eq:residual-main}
\end{equation}
with \(\ufull=\ubg+\ures\).  Equation~\eqref{eq:residual-main} is an exact
change of variables, not a linearization.  It retains nonlinear cross terms
between the background and the residual because \(\rhs[\ufull]\) is evaluated
on the full state.

For a large mass ratio this is especially useful: the primary background can
be large while the secondary field may be small over most of the domain.
Near the secondary, however, the residual may be order unity.  The method is
still internally consistent there; it is then best understood as a reference
metric formulation rather than a perturbative formulation.

\section{Gauge Variables}

The lapse and shift are split in the same way:
\begin{align}
  \afull &= \abg + \ares, \\
  \bfull &= \bbg + \bres.
\end{align}
The background gauge variables \(\abg,\bbg\) are the analytic lapse and shift
of the prescribed background.  They must not be modified locally by the
secondary object.

There are three useful gauge modes:
\begin{enumerate}[label=\arabic*.]
  \item \textbf{Pure prescribed background gauge.}
  Set \(\ares=0\) and \(\bres=0\), or preserve only selected initial residuals.
  This makes the full gauge equal to the background gauge except for any
  explicitly preserved residual components.

  \item \textbf{Background-subtracted dynamic residual gauge.}
  Evolve the residual lapse and shift using the difference between the full
  gauge RHS and the background gauge RHS.

  \item \textbf{Target-damped residual gauge.}
  Evolve a background-subtracted dynamic residual gauge, but add damping
  toward a prescribed residual target, usually for startup only.
\end{enumerate}

\section{Background-Subtracted Dynamic Gauge}

Let the ordinary full-spacetime gauge equations be written schematically as
\begin{align}
  \partial_t \alpha &= \gaugea[u,\alpha,\beta^i], \\
  \partial_t \beta^i &= \gaugeb[u,\alpha,\beta^i].
\end{align}
For example, \(\gaugea\) may represent a 1+log-like lapse condition and
\(\gaugeb\) may represent a Gamma-driver-like shift condition.

The residual gauge evolution should be
\begin{align}
  \partial_t \ares
  &=
  \gaugea[\ufull,\afull,\bfull]
  -
  \gaugea[\ubg,\abg,\bbg],
  \label{eq:alpha-residual-gauge}
  \\
  \partial_t \bres
  &=
  \gaugeb[\ufull,\afull,\bfull]
  -
  \gaugeb[\ubg,\abg,\bbg].
  \label{eq:beta-residual-gauge}
\end{align}
This is the gauge analogue of Eq.~\eqref{eq:residual-main}.  The pure
background state is an exact fixed point of the residual system: if
\(\ures=\ares=\bres=0\), then the residual RHS vanishes even if the analytic
background is not a stationary solution of the ordinary full gauge condition.

This point is essential.  A Kerr-Schild Schwarzschild background is not
generally stationary under a generic moving-puncture 1+log/Gamma-driver gauge.
If one evolved the full gauge equation directly on top of the background,
large gauge motion of the primary black hole would be injected into the
residual fields.  Equations~\eqref{eq:alpha-residual-gauge} and
\eqref{eq:beta-residual-gauge} subtract that background gauge motion.

\section{Why This Does Not Feed Back Into the Background}

The background is not a separately evolved physical component.  It is a fixed
reference field.  Local residual gauge evolution changes the full spacetime:
\begin{equation}
  (\ufull,\afull,\bfull)
  =
  (\ubg+\ures,\abg+\ares,\bbg+\bres).
\end{equation}
It does not change
\begin{equation}
  (\ubg,\abg,\bbg).
\end{equation}

Nonlinear feedback from the local gauge choice is still present, but it appears
in the residual variables through \(\rhs[\ufull]-\rhs[\ubg]\).  That is the
correct place for it.  What would be inconsistent is evaluating the background
subtraction using the full local gauge, for example
\begin{equation}
  \rhs[\ubg,\afull,\bfull],
\end{equation}
or damping the analytic background variables toward a star- or secondary-BH
gauge.  Those choices would force the reference background to obey a local
gauge prescription that is not its analytic solution.

The consistency rule is therefore:
\begin{quote}
The background callback always returns the pure analytic background in its
own analytic gauge.  All local gauge differences belong to the residual state.
\end{quote}

\section{Target-Damped Residual Gauge}

For robustness one may add a target damping term:
\begin{align}
  \partial_t \ares
  &=
  \gaugea[\ufull,\afull,\bfull]
  -
  \gaugea[\ubg,\abg,\bbg]
  -
  \eta_{\alpha}(t,x)\left(\ares-\ares^{\mathrm{tar}}\right),
  \label{eq:target-alpha}
  \\
  \partial_t \bres
  &=
  \gaugeb[\ufull,\afull,\bfull]
  -
  \gaugeb[\ubg,\abg,\bbg]
  -
  \eta_{\beta}(t,x)\left(\bres-\beta_{\mathrm{res}}^{i,\mathrm{tar}}\right).
  \label{eq:target-beta}
\end{align}
The target fields should be residual targets, not background modifications:
\begin{align}
  \ares^{\mathrm{tar}}(x,t) &=
  w(t,x)\left(\alpha_{\mathrm{TOV/boost}}(x,t)
  - \abg(x,t)\right),\\
  \beta_{\mathrm{res}}^{i,\mathrm{tar}}(x,t) &=
  w(t,x)\left(\beta^i_{\mathrm{TOV/boost}}(x,t)
  - \bbg(x,t)\right).
\end{align}
For a star, the target may be the analytic static or boosted TOV gauge used in
the initial data.  For a secondary black hole, a target should probably be used
only as a startup driver, because a permanently prescribed secondary-BH gauge
may fight the nonlinear two-body geometry.

A simple time ramp is
\begin{equation}
  w(t) = \exp\left[-\left(t/\tau\right)^p\right],
\end{equation}
with \(\tau \sim 50M\) and \(p=2\) or \(4\).  A compact spatial window around
the secondary can also be used if the target should act only where the
secondary gauge is meaningful.

The target-damped scheme remains a gauge prescription.  It is not a physical
source term, provided it vanishes for the pure background case and is applied
only to residual gauge variables.

\section{Relation to Current AthenaK Options}

The current code already contains the main ingredients:
\begin{itemize}
  \item \verb|use_analytic_background = true| selects the residual/background
  path.
  \item \verb|evolve_gauge_residual = false| prescribes gauge residual RHS
  to zero.  In that mode the full lapse and shift are the analytic background
  values unless residual components are preserved.
  \item \verb|preserve_lapse_residual = true| keeps the initial lapse residual
  while still prescribing its RHS to zero.  This is a static residual-gauge
  target for the lapse.
  \item \verb|evolve_gauge_residual = true| is intended to compute the full
  gauge RHS and subtract the background gauge RHS.
\end{itemize}

The debugging evidence so far suggests that pinning the full lapse to the
background value can be too crude for a self-gravitating TOV star.  In the
short static-Minkowski discriminator, preserving the initial TOV lapse residual
substantially reduced early metric drift: \(\det\gamma_{\min}\) stayed close
to unity and the growth of \(K\)-related diagnostics was strongly suppressed.

For a boosted star, the shift residual may also be relevant.  The pgen appears
to construct a boosted lab-frame lapse and shift for the initial star.  If
\verb|evolve_gauge_residual = false| and only the lapse residual is preserved,
then the boosted-star shift residual is still discarded.  A natural extension
is therefore a controlled \verb|preserve_shift_residual| or
\verb|target_shift_residual| option.

\section{Recommended Test Matrix}

For the current Schwarzschild radial-infall problem, the recommended order is:
\begin{enumerate}[label=\arabic*.]
  \item \textbf{Preserve lapse residual only.}
  Run with
  \begin{verbatim}
  <z4c>
  evolve_gauge_residual = false
  preserve_lapse_residual = true
  diss = 0.5  # or 0.6 as a stability-margin check
  damp_kappa1 = 0.0
  \end{verbatim}
  Compare the first non-finite time and the growth of Z4c constraints against
  the previous Schwarzschild run, which first became non-finite near \(t=80M\).

  \item \textbf{Pure dynamic residual gauge.}
  Run with
  \begin{verbatim}
  <z4c>
  evolve_gauge_residual = true
  preserve_lapse_residual = false
  \end{verbatim}
  This tests whether the background-subtracted 1+log/Gamma-driver path is
  already stable without target damping.

  \item \textbf{Target-damped lapse residual.}
  Add a damping term of the form Eq.~\eqref{eq:target-alpha}, with the initial
  TOV or boosted-TOV lapse residual as the target.  First try no spatial
  window, then a secondary-centered window if needed.

  \item \textbf{Target-damped shift residual.}
  If the boosted-star runs remain unstable, add the analogous shift target.
  This is especially relevant because a boosted star has a nonzero local
  shift in the lab-frame initial data.

  \item \textbf{Dissipation and damping margin checks.}
  Only after the gauge-path discriminator is clear, vary KO dissipation and
  constraint damping.  Dissipation may delay a gauge-driven crash, but it
  should not be treated as the primary fix if the underlying lapse/shift
  prescription is inconsistent with the self-gravitating secondary.
\end{enumerate}

\section{Diagnostics}

Each run should record:
\begin{itemize}
  \item first non-finite time in \verb|mhd.hst| and \verb|z4c.user.hst|;
  \item constraint norms \(C,H,M,Z,\Theta\) versus time;
  \item \(\rho_{\max}\), mass, and star tracker location versus time;
  \item \(\alpha_{\min}\), \(\alpha_{\max}\), and \(\max |\beta^i|\);
  \item metric health diagnostics such as \(\det\gamma_{\min}\),
  \(\chi_{\min}\), \(K\)-related maxima, and bad-metric counts;
  \item y-reflection and z-reflection symmetry norms on dense slices, with
  high-density masks.
\end{itemize}

A useful failure classification is:
\begin{quote}
If Z4c constraints grow rapidly while density remains finite, and MHD becomes
non-finite only after the Z4c diagnostics blow up, the leading failure mode is
residual spacetime/gauge instability rather than primary fluid dispersal.
\end{quote}

\section{Applicability to a Secondary Black Hole}

The residual formulation can be internally consistent for a secondary black
hole as well as for a star.  There is no formal small-residual assumption in
Eq.~\eqref{eq:residual-main}.  However, the practical demands are higher:
\begin{itemize}
  \item the secondary horizon scale must be resolved by AMR;
  \item the residual near the secondary can be order unity;
  \item the gauge must avoid lapse collapse, shift pathologies, and poor
  horizon treatment;
  \item the initial data may contain significant constraint violation if it is
  only a superposition rather than a solved binary-BH data set.
\end{itemize}

For a secondary black hole, a startup target for lapse and shift can help avoid
an abrupt gauge shock, but a permanent analytic target is risky.  A better
strategy is to ramp the target away and let a background-subtracted residual
moving-puncture gauge control the secondary.

\section{Implementation Checklist}

The safest implementation sequence is:
\begin{enumerate}[label=\arabic*.]
  \item Confirm that background state construction always writes pure analytic
  background lapse and shift.
  \item Confirm that full-state reconstruction uses
  \(\afull=\abg+\ares\) and \(\bfull=\bbg+\bres\).
  \item Confirm that gauge RHS subtraction evaluates
  \(\gaugea[\ufull]-\gaugea[\ubg]\) and
  \(\gaugeb[\ufull]-\gaugeb[\ubg]\) with the same derivative operators.
  \item Add optional residual target damping only to residual gauge RHS.
  \item Keep KO dissipation acting on residual fields, not on the reconstructed
  full state.
  \item Verify pure-background tests: with no star or secondary, residuals
  should remain at roundoff under the chosen gauge path.
  \item Verify static TOV on Minkowski before boosted Schwarzschild infall.
\end{enumerate}

\section{Summary}

Prescribing the analytic background lapse and shift is consistent only if they
remain properties of the fixed reference background.  Local gauge changes
required by the secondary must be carried by residual lapse and shift
variables.  A background-subtracted residual gauge evolution provides a
mathematically consistent way to let the secondary gauge adjust while keeping
the primary background subtraction clean.  Target damping of the residual
gauge can be used as a startup stabilizer, but it should be introduced as a
residual gauge prescription rather than as any modification of the analytic
background.

\end{document}
